\section{Characterizations of alpha equivalence}
\label{sec:alpha-equivalences}

The literature defines alpha equivalence of named terms in several
ways: as the congruence generated by an axiom for renaming a bound
atom, with one of two side conditions on the new atom, or, following
Gabbay and Pitts~\cite{gabbay2002a}, by a syntax-directed system of
rules in which the two bound atoms are swapped with a third. That
these definitions agree is folklore, but
Crole~\cite{crole2012a} observes that the proof sketches in the
literature are sometimes incorrect, and gives a detailed proof for
five such relations. This section presents his result as it is
formalized in CSLib. We keep Crole's names for the relations, the
rules and the side conditions; each statement names its counterpart
in~\cite{crole2012a}, and a raised \leanicon{} links to its Lean
source.

\subsection{Swapping}
\label{subsec:swapping}

\begin{definition}[Permutation action {\cite[Section~2]{crole2012a}}]
\label{def:perm-action}
\lean{\CroleSrc/Named/Untyped/AlphaEquivDefs.lean\#L68}
For a permutation $\pi$ of $\Atoms$, the term $\pi \act M$ is defined by
\[
  \pi \act a = \pi(a)
  \qquad
  \pi \act (M_1\, M_2) = (\pi \act M_1)\,(\pi \act M_2)
  \qquad
  \pi \act (\lam{a}{M}) = \lam{\pi(a)}{\pi \act M}.
\]
For the transposition $\sw{a}{b}$ we call $\sw{a}{b} \act M$ the
\emph{swap}
\lean{\CroleSrc/Named/Untyped/AlphaEquivDefs.lean\#L80}
of $a$ and $b$ in $M$. The \emph{agreement set}
\lean{\CroleSrc/Named/Untyped/AlphaEquivDefs.lean\#L60}
of two permutations is
$\AS(\pi, \pi') = \{\, a \in \Atoms \mid \pi(a) = \pi'(a) \,\}$.
\end{definition}

Unlike renaming (Definition~\ref{def:rename}), a swap is a bijection
on terms and never identifies two atoms. The two operations agree
when the new atom is unused: if $z \nocc M$ then
$\sw{z}{a} \act M = \ren{M}{a}{z}$
\lean{\CroleSrc/Named/Untyped/Properties.lean\#L159}.

\subsection{The five relations}
\label{subsec:five-definitions}

The first three relations are syntax-directed and differ only in the
side condition on the atom $z$ that replaces the two bound atoms.

\begin{definition}[{\cite[Definitions~3.1--3.3]{crole2012a}}]
\label{def:simp}
The relation $\simp$
\lean{\CroleSrc/Named/Untyped/Basic.lean\#L84}
on $\Lam$ is inductively defined by
\begin{gather*}
  \infrule{atom}{}{a \simp a}
  \qquad
  \infrule{pcg}{M_1 \simp M_1' \quad M_2 \simp M_2'}
               {M_1\, M_2 \simp M_1'\, M_2'}
  \qquad
  \infrule{pi}{\sw{z}{a} \act M \simp \sw{z}{b} \act M'}
              {\lam{a}{M} \simp \lam{b}{M'}}
  \; [z \nocc a, b, M, M'].
\end{gather*}
The relations $\simpf$
\lean{\CroleSrc/Named/Untyped/AlphaEquivDefs.lean\#L89}
and $\simpi{1}$
\lean{\CroleSrc/Named/Untyped/AlphaEquivDefs.lean\#L106}
are defined by the same rules, with the side condition of \textsc{pi}
replaced by
\[
  [z \fresh a, b, M, M']
  \qquad\text{and}\qquad
  [z \nocc M, M'],
\]
respectively; we call the resulting rules \textsc{pi\#} and
\textsc{pi1}.
\end{definition}

The relation $\simp$ follows Gabbay and Pitts~\cite{gabbay2002a} and
is CSLib's alpha equivalence $\alphaeq$. The remaining two relations
are generated by the renaming axiom and are congruences by
definition.

\begin{definition}[{\cite[Definitions~3.4 and~3.5]{crole2012a}}]
\label{def:simr}
The relation $\simr$
\lean{\CroleSrc/Named/Untyped/AlphaEquivDefs.lean\#L145}
on $\Lam$ is inductively defined by
\begin{gather*}
  \infrule{ref}{}{M \simr M}
  \qquad
  \infrule{sym}{M_1 \simr M_2}{M_2 \simr M_1}
  \qquad
  \infrule{trs}{M_1 \simr M_2 \quad M_2 \simr M_3}{M_1 \simr M_3}
  \\[1ex]
  \infrule{pcg}{M_1 \simr M_1' \quad M_2 \simr M_2'}
               {M_1\, M_2 \simr M_1'\, M_2'}
  \qquad
  \infrule{bcg}{M \simr M'}{\lam{a}{M} \simr \lam{a}{M'}}
  \qquad
  \infrule{$\alpha$}{}{\lam{a}{M} \simr \lam{a'}{\subst{M}{a}{a'}}}
  \; [a' \nocc a, M].
\end{gather*}
The relation $\simrf$
\lean{\CroleSrc/Named/Untyped/AlphaEquivDefs.lean\#L168}
is defined by the same rules, with the side condition of $\alpha$
replaced by $[a' \fresh a, M]$; we call the resulting rule
$\alpha\#$.
\end{definition}

The relation $\simr$ corresponds to the congruence $\equiv_\alpha$
of Barendregt~\cite[Definition~2.1.11]{barendregt1985a}, and $\simrf$
to that of Hindley and Seldin~\cite[Definition~1.17]{hindley2008a}.

\subsection{Properties of $\simp$}
\label{subsec:key-lemmas}

Since $\simp$ has no rules for reflexivity, symmetry and
transitivity, these have to be proved. They were already available in
CSLib, together with the two other basic facts below; the
corresponding statements for $\equiv_\alpha$
are~\cite[Lemmas~1.19 and~1.21]{hindley2008a}, proved
in~\cite[Appendix~A1]{hindley2008a}.

\begin{lemma}
\label{lem:equivalence}
The relation $\simp$ is
reflexive~\lean{\CroleSrc/Named/Untyped/Properties.lean\#L493},
symmetric~\lean{\CroleSrc/Named/Untyped/Properties.lean\#L373}
and
transitive~\lean{\CroleSrc/Named/Untyped/Properties.lean\#L566}.
Moreover, if $M \simp M'$ and $N \simp N'$ then
$\free(M) = \free(M')$
\lean{\CroleSrc/Named/Untyped/Properties.lean\#L122}
and $\subst{M}{a}{N} \simp \subst{M'}{a}{N'}$
\lean{\CroleSrc/Named/Untyped/Properties.lean\#L902}.
\end{lemma}

The next two lemmas say that $\simp$ is invariant under swapping,
an instance of the equivariance principle of nominal
techniques~\cite[Theorem~4.4]{gabbay2011a}, and that the action of a
permutation on a term depends, up to $\simp$, only on its values on
the free atoms.

\begin{lemma}[{\cite[Lemma~6.1]{crole2012a}}]
\label{lem:swap-preserve}
\lean{\CroleSrc/Named/Untyped/Properties.lean\#L381}
If $M \simp M'$ then $\sw{u}{v} \act M \simp \sw{u}{v} \act M'$, for
all atoms $u, v$.
\end{lemma}

\begin{lemma}[{\cite[Lemma~6.2]{crole2012a}}]
\label{lem:agreement}
Let $\pi, \pi'$ be permutations of $\Atoms$ and $M$ a term.
\begin{enumerate}
  \item \lean{\CroleSrc/Named/Untyped/AlphaEquivLemmas.lean\#L32}
    If $\occ(M) \subseteq \AS(\pi, \pi')$ then
    $\pi \act M = \pi' \act M$.
  \item \lean{\CroleSrc/Named/Untyped/AlphaEquivLemmas.lean\#L63}
    If $\free(M) \subseteq \AS(\pi, \pi')$ then
    $\pi \act M \simp \pi' \act M$.
\end{enumerate}
\end{lemma}

The bridge between the two families of relations is that swapping
with an unused atom is, up to alpha equivalence, substitution. For
the renaming relations this
is~\cite[Lemmas~6.3 and~6.4]{crole2012a}. Gabbay proves it as a
syntactic identity under the additional hypothesis that $a$ is not
bound in $M$~\cite[Lemma~5.16]{gabbay2011a}.

\begin{lemma}
\label{lem:swap-subst}
If $z \nocc M$ then
$\sw{z}{a} \act M \simr \subst{M}{a}{z}$
\lean{\CroleSrc/Named/Untyped/AlphaEquivLemmas.lean\#L167}
and
$\sw{z}{a} \act M \simrf \subst{M}{a}{z}$
\lean{\CroleSrc/Named/Untyped/AlphaEquivLemmas.lean\#L203}.
\end{lemma}

For $\simp$ the corresponding statement compares two swaps.

\begin{lemma}[{\cite[Lemma~6.5]{crole2012a}}]
\label{lem:swap-chain}
\lean{\CroleSrc/Named/Untyped/AlphaEquivLemmas.lean\#L402}
Let $z_1, a_1, a_1', \dots, z_n, a_n, a_n'$ be pairwise distinct atoms
and let $M$ be a term with $z_i \fresh M$ and $a_i' \fresh M$ for all
$i$. Then
\[
  \sw{z_1}{a_1} \act \dots \act \sw{z_n}{a_n} \act M
  \;\simp\;
  \sw{z_1}{a_1'} \act \dots \act \sw{z_n}{a_n'} \act
    \subst{\subst{M}{a_1}{a_1'} \dots}{a_n}{a_n'}.
\]
\end{lemma}

Only the case $n = 1$
\lean{\CroleSrc/Named/Untyped/AlphaEquivLemmas.lean\#L476}
is used below. Crole assumes $z_i \nocc M$; freshness suffices,
because a bound occurrence of $z_i$ can be renamed away first:

\begin{lemma}[{\cite[Lemma~6.6]{crole2012a}}]
\label{lem:avoid}
If $a \fresh M$ then there is a term $\hat{M}$ with $a \nocc \hat{M}$
and $\hat{M} \simr M$
\lean{\CroleSrc/Named/Untyped/AlphaEquivLemmas.lean\#L210},
and likewise with $\simp$ in place of $\simr$
\lean{\CroleSrc/Named/Untyped/AlphaEquivLemmas.lean\#L243}.
\end{lemma}

This lemma justifies Barendregt's variable
convention~\cite[2.1.13]{barendregt1985a}; for $\equiv_\alpha$ it
is~\cite[Lemma~A1.5(g)]{hindley2008a}. The last lemma is the
counterpart of the second half of Lemma~\ref{lem:equivalence} for the
renaming relation, and a special case
of~\cite[Proposition~C.2]{barendregt1985a}
and~\cite[Lemma~A1.10]{hindley2008a}.

\begin{lemma}[{\cite[Lemma~6.7]{crole2012a}}]
\label{lem:subst-congr}
\lean{TODO}
If $M \simr M'$ then $\subst{M}{b}{a} \simr \subst{M'}{b}{a}$, for
all atoms $a, b$.
\end{lemma}
%% TODO(Chris): no Lean declaration for this lemma (for \simr) exists at the pinned commit.

\subsection{The characterization theorem}
\label{subsec:equivalence-theorem}

\begin{theorem}[Crole {\cite[Theorems~4.1, 4.2, 4.4--4.6]{crole2012a}}]
\label{thm:crole-equivalence}
For all terms $M, N \in \Lam$ the following are equivalent:
\begin{enumerate}
  \item[(i)] $M \simp N$;
  \item[(ii)] $M \simpf N$
    \lean{\CroleSrc/Named/Untyped/AlphaEquivEquiv.lean\#L102};
  \item[(iii)] $M \simpi{1} N$
    \lean{\CroleSrc/Named/Untyped/AlphaEquivEquiv.lean\#L151};
  \item[(iv)] $M \simr N$
    \lean{\CroleSrc/Named/Untyped/AlphaEquivEquiv.lean\#L261};
  \item[(v)] $M \simrf N$
    \lean{\CroleSrc/Named/Untyped/AlphaEquivEquiv.lean\#L324}.
\end{enumerate}
\end{theorem}

\begin{proof}
Each of (ii)--(v) is shown equivalent to (i), by showing that each of
the two relations is closed under the rules of the other. Only the
rules for abstractions need an argument.

\emph{(i)\,$\Leftrightarrow$\,(ii).} Non-occurrence implies
freshness, so \textsc{pi} is an instance of \textsc{pi\#}.
Conversely, let $\sw{u}{a} \act M \simp \sw{u}{b} \act M'$ with
$u \fresh a, b, M, M'$, and choose $z \nocc u, a, b, M, M'$. By
Lemma~\ref{lem:swap-preserve},
$\sw{z}{u} \act \sw{u}{a} \act M \simp \sw{z}{u} \act \sw{u}{b} \act M'$.
The permutations $\sw{z}{u} \circ \sw{u}{a}$ and $\sw{z}{a}$ agree on
$\free(M)$, so the left-hand side is $\simp$-related to
$\sw{z}{a} \act M$ by Lemma~\ref{lem:agreement}, and likewise for the
right-hand side. Hence $\sw{z}{a} \act M \simp \sw{z}{b} \act M'$, and
\textsc{pi} applies.

\emph{(i)\,$\Leftrightarrow$\,(iii).} The rule \textsc{pi1} is more
general than \textsc{pi} only in allowing $z \in \{a, b\}$. If
$z = a$, its premise reads $M \simp \sw{a}{b} \act M'$ with
$a \nocc M'$. For $z' \nocc a, b, M, M'$,
Lemma~\ref{lem:swap-preserve} and the first part of
Lemma~\ref{lem:agreement} give
$\sw{z'}{a} \act M \simp \sw{z'}{b} \act M'$. The case $z = b$ is
symmetric.

\emph{(i)\,$\Leftrightarrow$\,(iv).} For \textsc{pi}, let
$\sw{z}{a} \act M \simr \sw{z}{b} \act M'$ with
$z \nocc a, b, M, M'$. Lemma~\ref{lem:swap-subst} gives
$\subst{M}{a}{z} \simr \subst{M'}{b}{z}$, and then
\[
  \lam{a}{M} \;\simr\; \lam{z}{\subst{M}{a}{z}}
  \;\simr\; \lam{z}{\subst{M'}{b}{z}} \;\simr\; \lam{b}{M'}
\]
by $\alpha$, \textsc{bcg} and $\alpha$. Conversely, \textsc{ref},
\textsc{sym} and \textsc{trs} hold for $\simp$ by
Lemma~\ref{lem:equivalence}, and \textsc{bcg} by
Lemma~\ref{lem:swap-preserve} and \textsc{pi}. For $\alpha$, let
$a' \nocc a, M$ and choose $z \nocc a, a', M, \subst{M}{a}{a'}$. Then
$\sw{z}{a} \act M \simp \sw{z}{a'} \act \subst{M}{a}{a'}$ by
Lemma~\ref{lem:swap-chain}, which is the premise of \textsc{pi} for
$\lam{a}{M} \simp \lam{a'}{\subst{M}{a}{a'}}$.

\emph{(i)\,$\Leftrightarrow$\,(v).} As in the previous case:
Lemma~\ref{lem:swap-subst} covers $\simrf$, and
Lemma~\ref{lem:swap-chain} only requires $a' \fresh M$.
\end{proof}

In the last case Crole's argument differs from ours. He proves the
rule $\alpha$ by an induction on $M$ that fails for $\alpha\#$, where
the substitution $\subst{M}{a}{a'}$ may have to rename a binder, and
introduces Lemma~\ref{lem:swap-chain} for general $n$ to repair it.
Proving the case $n = 1$ under the freshness hypothesis serves both
rules.

The equivalence of (iv) and (v) follows from the theorem, but it is a
statement about renaming alone, and Crole also proves it without
recourse to
swapping~\cite[Theorem~4.6]{crole2012a}
\lean{\CroleSrc/Named/Untyped/AlphaEquivEquiv.lean\#L337}.
%% TODO(Chris): this declaration (alphaEquivR_iff_alphaEquivRFresh_direct) is still `sorry`
%% at the pinned commit. Finish it and re-pin before submission.
Since non-occurrence implies freshness, $\alpha$ is an instance of
$\alpha\#$. Conversely, let $a' \fresh a, M$. By
Lemma~\ref{lem:avoid} there is $\hat{M} \simr M$ with
$a' \nocc \hat{M}$, and
\[
  \lam{a}{M} \;\simr\; \lam{a}{\hat{M}}
  \;\simr\; \lam{a'}{\subst{\hat{M}}{a}{a'}}
  \;\simr\; \lam{a'}{\subst{M}{a}{a'}}
\]
by \textsc{bcg}, by $\alpha$, and by Lemma~\ref{lem:subst-congr} with
\textsc{bcg}. Hindley and Seldin relate $\equiv_\alpha$ to a
restricted renaming step by the same
argument~\cite[Lemmas~A1.7 and~A1.8]{hindley2008a}.

The side condition of \textsc{pi} cannot be weakened much further.

\begin{proposition}[{\cite[Proposition~4.3]{crole2012a}}]
\label{prop:variants}
\lean{\CroleSrc/Named/Untyped/AlphaEquivEquiv.lean\#L157}
Let $\simpi{2}$
\lean{\CroleSrc/Named/Untyped/AlphaEquivDefs.lean\#L118}
and $\simpi{3}$
\lean{\CroleSrc/Named/Untyped/AlphaEquivDefs.lean\#L129}
be defined as $\simp$, with the side condition of \textsc{pi} replaced
by $[z \nocc a, b]$ and by no condition, respectively. Then
${\simpi{2}} \neq {\simp}$ and ${\simpi{3}} \neq {\simp}$.
\end{proposition}

\begin{proof}
For distinct atoms $a, b, z$ we have
$\lam{a}{a\, z} \simpi{2} \lam{b}{b\, a}$ and
$\lam{a}{b\, a} \simpi{3} \lam{b}{a\, b}$, but in both cases the two
terms have different free atoms, so they are not related by $\simp$
(Lemma~\ref{lem:equivalence}).
\end{proof}

\subsection{The formalization}
\label{subsec:formalization-notes}

The formalization follows~\cite{crole2012a} closely, and we found no
errors in it. To our knowledge the result had not been mechanized
before; Crole mentions an Isabelle/HOL development in
progress~\cite[Section~7]{crole2012a}, of which we could find no
published account.
%% TODO(Chris): confirm the novelty claim against the literature before submission. Related
%% but different: Apinis and Ahman, "A simple formalization of alpha-equivalence"
%% (arXiv:2507.10181, Rocq), relate one inductive definition to the literature's.
Our contribution is the four relations other than $\simp$,
Theorem~\ref{thm:crole-equivalence},
Proposition~\ref{prop:variants}, and
Lemmas~\ref{lem:swap-preserve}--\ref{lem:subst-congr}. The definition of
$\simp$ and Lemma~\ref{lem:equivalence} belong to the named
representation that this work builds on
(Section~\ref{subsec:named-terms}). The first part of the development
is under review for inclusion in
CSLib\footnote{\url{https://github.com/leanprover/cslib/pull/741}};
the remainder is to follow in the same way.

What a formalization has to add to~\cite{crole2012a} are the proofs
that the paper leaves to the reader or to the literature.
Table~\ref{tab:gaps} lists them, with a place where the argument can
be consulted on paper. None of them held a surprise.

\begin{table}[b]
\caption{Proofs omitted or only sketched in~\cite{crole2012a}.}
\label{tab:gaps}
\begin{tabular}{@{}lll@{}}
\toprule
Statement in~\cite{crole2012a} & Treatment there & See \\
\midrule
Existence of substitution (Theorem~A.1) & deferred
  & \cite[Definition~5.13, Lemma~5.14]{gabbay2011a} \\
Lemma~6.1 & representative cases only
  & Lemma~\ref{lem:swap-preserve}; \cite[Theorem~4.4]{gabbay2011a} \\
Lemma~6.2, part~1 & omitted, as for part~2
  & Lemma~\ref{lem:agreement} \\
Lemma~6.5 & details omitted & Lemma~\ref{lem:swap-chain} \\
Lemma~6.6 & sketched as well known
  & \cite[Lemma~A1.5(g)]{hindley2008a} \\
Lemma~6.7 & deferred as well known
  & \cite[Lemma~A1.10]{hindley2008a}; \cite[pp.~96--103]{curry1958a} \\
\bottomrule
\end{tabular}
\end{table}
%% TODO(Chris): the Curry--Feys pages are quoted from the proof of Hindley--Seldin's
%% Lemma A1.10(a) ("[CF58, Theorem 2(a) p. 95, proof on pp. 96-103]"), not read first-hand.
%% gabbay2011a numbers follow the 54-page author version, which is also what Crole cites.

Three statements differ from their counterparts
in~\cite{crole2012a}. The rules \textsc{pi} and \textsc{pi1} are
stated in CSLib with renaming (Definition~\ref{def:rename}) in place
of the swap, which is the same under their side conditions.
Lemma~\ref{lem:agreement} is proved for arbitrary bijections on
$\Atoms$, where Crole assumes finite support; the difference is
immaterial for terms~\cite[Remark~2.25]{gabbay2011a}. And
Lemma~\ref{lem:swap-chain} has the weaker hypothesis discussed above.
Crole's second contribution, alpha equivalence of program
contexts~\cite[Section~5]{crole2012a}, is not covered.
